\chapter{Apr.~15 --- Quantum Groups, Part 3}

\section{Quantum Double for \texorpdfstring{$U_q(\g_e)$}{Uq(ge)}}

\begin{remark}
  We construct the quantum double for
  $U_q(\g_e)$.
  First consider
  $U_q(\g^{\pm}_e)$, generated by
  $\{q^h, e_i\}$
  and $\{q^h, f_i\}$, respectively:
  \[
    U_q(\g^{\pm}_e)
    = \C[\mathfrak{h}_e]
    \otimes U_q(\mathfrak{n}^{\pm}),
  \]
  where $\C[\mathfrak{h}_e]$ is the group algebra
  associated with $q^h$. Then we can
  define the restricted dual
  (by degree)
  \[
    U_q(\g^{\pm})^*
    = \C[\mathfrak{h}_e^*] \otimes U_q(\mathfrak{n}^{\pm})^*,
  \]
  where the pairing is
  $\langle q^h, q^{\lambda} \rangle = q^{\lambda(h)}$
  for $\lambda \in \mathfrak{h}_e^*$.
\end{remark}

\begin{theorem}\label{thm:double}
  There is an isomorphism of Hopf algebras
  \[
    \theta :
    U_q(\g_e^-)
    \longrightarrow
    (U_q(\g_e^+)^*)^{\mathrm{op}}
  \]
  such that
  $\theta(q^h) = q^h$
  for all $h \in \mathfrak{h}_e$,
  where we identify
  $\mathfrak{h}_e$ with $\mathfrak{h}_e^*$
  using $\llparenthesis \cdot, \cdot \rrparenthesis$.
\end{theorem}

\begin{example}
  Check that Theorem \ref{thm:double}
  does not hold for $q = 1$ (one algebra
  is commutative whereas the other is not).
\end{example}

\begin{theorem}
  Theorem \ref{thm:double} implies that
  we can identify
  \[
    D(U_q(\g_e^+))
    = U_q(\g_e^+) \otimes U_q(\g_e^-),
  \]
  which has a two-sided ideal $H$ generated
  by $q^h \otimes 1 - 1 \otimes q^h$.
\end{theorem}

\begin{theorem}
  The map given by
  \[
    e_i \otimes 1 \mapsto e_i, \quad
    1 \otimes f_i \mapsto f_i, \quad
    q^h \otimes 1 \mapsto q^h, \quad
    1 \otimes q^h \mapsto q^{h}
  \]
  defines an isomorphism
  $\phi : (U_q(\g_e^+) \otimes U_q(\g_e^-)) / H \to U_q(\g_e)$.
\end{theorem}

\begin{remark}
  The algebra $D(U_q(\g_e^+))$ has
  $R$-matrix given by
  \[
    R = q^{\sum_p x_p \otimes x_p}
    \sum_i a_i \otimes a^i
    = \sum_{n \ge 0}
    \frac{t^n (\sum_p x_p \otimes x_p)}{n!}
    \sum_i a_i \otimes a^i,
  \]
  where $\{a_i\}$ is a basis in
  $U_q(\mathfrak{n}^+)$ and
  $\{a^i\}$ is a basis in $U_q(\mathfrak{n}^-)$.
\end{remark}

\begin{example}
  For $U_q(\mathfrak{sl}_2)$, we have
  \[
    R = q^{\frac{1}{2} h \otimes h}
    \sum_{n \ge 0} q^{n(n - 1) / 2}
    \frac{(q - q^{-1}) e^n \otimes f^n}{[n]!},
  \]
  where $[n] = (q^n - q^{-n}) / (q - q^{-1})$.
  For general $\g_e$, the formula is given by
  Khoroshkin-Tolstoy.
\end{example}

\pagebreak

\begin{theorem}
  Let $C(\g_e, q)$ be the category of
  finite-dimensional representations
  with weight decomposition
  of $U_q(\g_e)$ (with $q$ not a root of
  unity). Then
  $C(\g_e, q)$ has the structure of
  a braided tensor category with
  $\sigma_{V, W} = \widecheck{R}_{V, W} = P_{1, 2} R_{V, W}$
  and unity associativity morphism.
\end{theorem}

\section{Quantum Casimir Element}

\begin{remark}
  Let $A$ be a quasitriangular Hopf algebra,
  and define
  \[
    u = m((\gamma \otimes {\id_A})(R^{\mathrm{op}}))
    = \sum_i \gamma(b_i) a_i,
  \]
  where $R = \sum_i a_i \otimes b_i$. Then
  $u \in A$ is well-defined
  (i.e. $u$ actually lies in $A$)
  if $A$ is finite-dimensional.
  In the case where $A = U_q(\g_e)$,
  the element is still defined as an
  operator in any
  finite-dimensional representation.
\end{remark}

\begin{prop}
  We have the following:
  \begin{enumerate}
    \item $\gamma^2(x) = u x u^{-1}$ for
      all $x \in A$;
    \item $u^{-1} = m((\gamma^{-1} \otimes {\id_A})(R^{\mathrm{op}}))$;
    \item $\Delta(u) = (u \otimes u)(R^{\mathrm{op}} R)^{-1}
      = (R^{\mathrm{op}} R)^{-1} (u \otimes u)$.
  \end{enumerate}
\end{prop}

\begin{remark}
  Recall that we have
  defined $q^{2\rho} \in U_q(\g_e)$
  satisfying
  $\gamma^2(x) = q^{2\rho} x q^{-2\rho}$.
\end{remark}

\begin{theorem}\label{thm:casimir}
  Let $A = U_q(\g_e)$ and
  $C = q^{2\rho} u^{-1}$. Then
  $C$ is a central element
  and is well-defined for any finite-dimensional
  representation. Moreover, $C$ satisfies
  \[
    \Delta(C)
    = (R^{\mathrm{op}} R) (C \otimes C)
    = (C \otimes C) (R^{\mathrm{op}} R)
  \]
  and
  $C|_V = q^{\llparenthesis \lambda, \lambda + 2\rho \rrparenthesis} \id_V$,
  where $V$ is the highest weight module
  with highest weight $\lambda$.
\end{theorem}

\begin{corollary}
  Let $V_\lambda, V_\mu, V_\nu$ be
  highest weight modules over
  $U_q(\g_e)$ with
  corresponding highest weights $\lambda, \mu, \nu$,
  respectively. Then for any
  embedding of $U_q(\g_e)$-modules
  $V_\lambda \subseteq V_\mu \otimes V_\nu$,
  we have
  \[
    \widecheck{R}^2|_{V_\lambda \subseteq V_\mu \otimes V_\nu}
    = q^{c(\lambda) - c(\mu) - c(\nu)} \id_{V_\lambda},
  \]
  where $c(\lambda) = \llparenthesis \lambda, \lambda + 2\rho \rrparenthesis$.
\end{corollary}

\begin{proof}
  Apply $\widecheck{R}^2 = \Delta(C)(C^{-1} \otimes C^{-1})$
  (which follows from Theorem \ref{thm:casimir})
  to $C|_V = q^{\llparenthesis \lambda, \lambda + 2\rho \rrparenthesis} \id_V$.
\end{proof}

\section{Intertwining Operators}

\begin{remark}
  Let $\g$ be a finite-dimensional simple Lie
  algebra
  and $C(\g, q)$ the category of
  finite-dimensional representations of
  $U_q(\g)$ with weight decomposition. Then
  all
  morphisms in this category
  can be described using
  $\Hom_{U_q(\g)}(L^{(q)}_\lambda, V)$.
  Let
  \[
    H_{V, \mu}^\lambda
    = \Hom_{U_q(\g)}(L^{(q)}_\lambda, V \otimes L^{(q)}_\mu).
  \]
  Let $V, W$ be finite-dimensional
  representations. Then
  \[
    \bigoplus_{\mu}
    H_{V, \mu}^\lambda \otimes H_{W, \nu}^\mu
    \cong \Hom_{U_q(\g)}(L^{(q)}_\lambda, V \otimes W \otimes L^{(q)}_\nu)
  \]
  via the composition
  $\Phi_1 \otimes \Phi_2 \mapsto (1 \otimes \Phi_2) \circ \Phi_1$.
\end{remark}

\begin{remark}
  Let $\widecheck{R}^+ = PR$
  and $\widecheck{R}^- = P(R^{\mathrm{op}})^{-1} = (\widehat{R}^+)^{-1}$.
  We can act on an intertwiner
  \[
    \Phi
    : L^{(q)}_\lambda
    \longrightarrow V \otimes W \otimes L^{(q)}_\nu
  \]
  by multiplying by 
  $\widecheck{R}^{\pm}_{V, W} \otimes {\id_{L^{(q)}_\nu}}$,
  which gives another intertwiner
  $L^{(q)}_\lambda \to W \otimes V \otimes L^{(q)}_\nu$.
  This gives
  \[
    \widecheck{B}^{\pm}_{W, V}
    (\lambda, \nu)
    : \bigoplus_{\mu} H_{V, \mu}^\lambda \otimes H_{W, \nu}^\mu
    \overset{\cong}\longrightarrow
    \bigoplus_{\mu'} H_{W, \mu'}^\lambda \otimes H_{V, \nu}^{\mu'}.
  \]
  We call such
  $\widecheck{B}^{\pm}$
  \emph{exchange matrices}.
\end{remark}

\begin{prop}
  The exchange matrices satisfy
  \begin{enumerate}
    \item $\widecheck{B}^{\pm}_{1, 2} \widecheck{B}^{\pm}_{2, 3} \widecheck{B}^{\pm}_{1, 2} = \widecheck{B}^{\pm}_{2, 3} \widecheck{B}^{\pm}_{1, 2} \widecheck{B}^{\pm}_{2, 3}$,
      where both sides are viewed as
      isomorphisms
      \[\Hom(L^{(q)}_\lambda, V \otimes W \otimes U \otimes L^{(q)}_\nu) \to \Hom(L^{(q)}_\lambda, U \otimes W \otimes V \otimes L^{(q)}_\nu);\]
    \item $\widecheck{B}^{\pm} \widecheck{B}^{\mp} = \id$.
  \end{enumerate}
\end{prop}

\begin{remark}
  Write the associativity morphism as
  $\alpha : (V \otimes L^{(q)}_\nu) \otimes W \cong V \otimes (L^{(q)}_\nu \otimes W)$.
  Then we have
  \[
    \alpha_{V, W}(\lambda, \mu)
    : \bigoplus_{\mu} H_{V, \nu}^\mu
    \otimes H_{W, \mu}^\lambda
    \overset{\cong}{\longrightarrow}
    \bigoplus_{\mu'} H_{V, \mu'}^{\lambda}
    \otimes H_{W, \nu}^{\mu'}.
  \]
  For $\mathfrak{sl}_2$ and
  $U_q(\mathfrak{sl}_2)$, let
  $V, W$ be highest weight modules
  with highest weights $i, j$.
  Then $\alpha_{V, W}$ depends on
  $i, j, \lambda, \nu, \mu, \mu'$, and the
  elements are called the
  \emph{quantum $6j$-symbols}
  (or \emph{Racah coefficients}).
\end{remark}

\begin{remark}
  The hexagon axiom gives the diagram
  \[
    \begin{tikzcd}
      V \otimes (L^{(q)}_\nu \otimes W)
      \ar[r, "\alpha^{-1}"] \ar[d, "{\id} \otimes \widecheck{R}_{L^{(q)}_\nu, W}", swap]
       &
       (V \otimes L^{(q)}_\nu) \otimes W \ar[d, "\widecheck{R}_{V \otimes L^{(q)}_\nu, W}"] \\
      V \otimes (W \otimes L^{(q)}_\nu)
      \ar[r, "\beta^+_{L^{(q)}_\nu, V, W}"]
      & W \otimes (V \otimes L^{(q)}_\nu)
    \end{tikzcd}
  \]
  Thus we can write
  $\widecheck{B}_{V, W}(\lambda, \nu) = \widecheck{R}_2 \alpha_{V, W}^{-1}(\lambda, \nu) \widecheck{R}_1^{-1}$,
  where
  \[
    \widecheck{R}_1 : \bigoplus_{\mu'}
    H_{W, \nu}^{\mu'}
    \longrightarrow
    \bigoplus_\mu H^\mu_{V, \nu}
    \quad \text{and} \quad
    \widecheck{R}_2 : \bigoplus_\mu H^\lambda_{W, \mu}
    \longrightarrow
    \bigoplus_{\mu'} H_{W, \mu'}^{\lambda}.
  \]
  For $U_q(\mathfrak{sl}_2)$, we have
  $\widecheck{B} = D_2 \widecheck{\alpha} D_1$,
  where $D_1, D_2$ are diagonal.
\end{remark}

\begin{remark}
  Let $\lambda, \mu, \nu$ be generic, so
  that $L^{(q)}_\lambda \cong M^{(q)}_\lambda$,
  and $V, W$ finite-dimensional. Then
  we can describe
  $H^{\lambda}_{V, \mu}$ explicitly using
  the map
  \begin{align*}
    \langle \cdot \rangle
    : \Hom_{U_q(\g)}(M^{(q)}_\lambda, M^{(q)}_\mu \otimes V)
    &\longrightarrow V^{\lambda - \mu} \\
    \Phi
    &\longmapsto
    \langle \Phi \rangle
    = (v_\mu^* \otimes {\id_V})(\Phi v_\lambda),
  \end{align*}
  where $v_\lambda$ is a highest weight
  vector in $M^{(q)}_\lambda$ and
  $v_\mu^*$ is a lowest weight vector in
  $(M^{(q)}_\mu)^*$. This gives
  \[
    H^\lambda_{V, \mu}
    = \Hom_{U_q(\g)}(M^{(q)}_\lambda, V \otimes M^{(q)}_\mu)
    \cong V^{\lambda - \mu}.
  \]
  Thus for any generic $\lambda$
  and $v \in V^{\lambda - \mu}$, there
  is a unique operator
  \[
    \Phi^v_\lambda
    : M^{(q)}_\lambda
    \longrightarrow
    M^{(q)}_\mu \otimes V
  \]
  such that $\Phi^v_\lambda(v_\lambda) = v \otimes v_\mu + \text{lower order terms}$.
  Similarly, if $v \in V$ and $w \in W$
  are homogeneous vectors, we can define
  \[
    \Phi^{v \otimes w}_\lambda
    = ({\id} \otimes \Phi^w) \circ \Phi^v_\lambda
    : L^{(q)}_\lambda
    \longrightarrow
    V \otimes W \otimes L^{(q)}_\nu,
  \]
  where $\nu = \lambda + \mathrm{weight}(v) + \mathrm{weight}(w)$.
  Then we have
  \[
    (\widecheck{R}^{\pm}_{V, W} \otimes {\id})
    \Phi^{v_i \otimes w_j}_\lambda
    = \sum_{i, j}
    \Phi^{w_{j'} \otimes v_{i'}}_\lambda
    B^{i, j}_{i', j'}(\lambda)^{\pm}.
  \]
  Using this, we can define
  $\widecheck{B}^{\pm}_{V, W}(\lambda, \nu) : (V \otimes W)^{\lambda - \nu} \cong (W \otimes V)^{\lambda - \nu}$
  between weight subspaces.
  Summing over $\nu$, we get an operator
  \[
    B^{\pm}_{V, W}(\lambda)
    : V \otimes W \longrightarrow W \otimes V
  \]
  which preserves weights. The
  $B^{\pm}_{V, W}(\lambda)$ satisfy
  \emph{dynamical Yang-Baxter equations}.
\end{remark}
